\documentclass{amsart}
\usepackage{amsmath,amssymb,amsthm,hyperref}
\usepackage{algorithm}
\usepackage{algpseudocode}

\newtheorem{theorem}{Theorem}
\newtheorem{lemma}{Lemma}
\newtheorem{proposition}{Proposition}
\theoremstyle{definition}
\newtheorem{definition}{Definition}
\newtheorem{remark}{Remark}

\DeclareMathOperator{\supp}{supp}
\DeclareMathOperator{\err}{err}
\DeclareMathOperator*{\OR}{OR}
\newcommand{\eps}{\varepsilon}
\newcommand{\X}{\{0,1\}^n}

\begin{document}

\title{Efficient semi-supervised PAC learning of two-sided disjunctions}
\maketitle

\section{Statement and answer}

We work in the model described in the problem. Examples lie in $X=\X$. A
\emph{two-sided disjunction} is a pair $h=(h^+,h^-)$ where, for some subsets
$S^+,S^-\subseteq\{1,\dots,n\}$,
\[
  h^+(x)=\OR_{i\in S^+} x_i=\mathbf 1[\,S^+\cap P(x)\neq\emptyset\,],\qquad
  h^-(x)=\OR_{i\in S^-} x_i=\mathbf 1[\,S^-\cap P(x)\neq\emptyset\,],
\]
with $P(x):=\{i:x_i=1\}$ (and the empty disjunction is identically $0$). The
\emph{classifier} associated with $h$ uses $h^+$ only: $x$ is labeled $+1$ if
$h^+(x)=1$ and $-1$ otherwise. For a distribution $D$ on $X$ and a target
$h^\ast=(h^{\ast+},h^{\ast-})$, the error of a hypothesis classifier $g$ is
\[
  \err_{D,h^\ast}(g):=\Pr_{x\sim D}\bigl[g(x)\neq h^{\ast+}(x)\bigr].
\]
A hypothesis $h$ is \emph{compatible} with $D$ if for every $x\in\supp(D)$ we
have $h^+(x)=\lnot h^-(x)$, i.e.\ exactly one of $h^+(x),h^-(x)$ equals $1$. Let
$C$ be the class of two-sided disjunctions with this compatibility notion, and
for a fixed $D$ let
\[
  C_D:=\{h\in C: h \text{ is compatible with } D\}
\]
be the compatibility-filtered subclass. Throughout we assume
$\eps,\delta\in(0,\tfrac12)$, the standard PAC regime (any
$\delta\in(0,1)$ is reduced to this regime by the usual confidence
amplification, which only multiplies the labeled sample by an
$O(\log(1/\delta))$ factor).

\begin{definition}[Filtered sample complexity]\label{def:filt}
$m_{\mathrm{filt}}(D,\eps,\delta)$ is the smallest integer $m$ for which there
exists a (possibly randomized) learning rule that takes $m$ labeled examples
drawn i.i.d.\ from $D$ and labeled by an arbitrary target $h^\ast\in C_D$, and
outputs a classifier $g$ such that, \emph{for every} $h^\ast\in C_D$,
$\Pr[\err_{D,h^\ast}(g)\le\eps]\ge1-\delta$, the probability over the sample and
the learner's internal randomness.
\end{definition}

\begin{theorem}\label{thm:main}
The answer to the problem is \textbf{yes}. There is an algorithm $\mathcal A$
for the semi-supervised PAC model of $C$ such that, on every unknown
distribution $D$ with $C_D\neq\emptyset$ and every target $h^\ast\in C_D$:
\begin{enumerate}
\item[(a)] $\mathcal A$ runs in time polynomial in $n$, $1/\eps$ and
$1/\delta$;
\item[(b)] with probability at least $1-\delta$ the output classifier has error
at most $\eps$, and the number of labeled examples used is
\[
  m \;=\; \Bigl\lceil \tfrac{1}{\eps}\bigl(n\ln 2+\ln(1/\delta)\bigr)\Bigr\rceil
  \;\le\; p\bigl(n,1/\eps,\log(1/\delta)\bigr)\cdot
          m_{\mathrm{filt}}(D,\eps,\delta)
\]
for the fixed polynomial $p(n,1/\eps,\log(1/\delta))=\tfrac1\eps\bigl(n+\log_2(1/\delta)+1\bigr)$;
\end{enumerate}
while using a number of unlabeled examples that is polynomial in $n$, $1/\eps$
and $\log(1/\delta)$ (indeed, $\mathcal A$ needs no unlabeled examples).
\end{theorem}

The bound in (b) is the genuine \emph{multiplicative} comparison requested in
the problem: there is no additive slack. The crux of the proof, and the
correction relative to a naive ``unfiltered $\le$ poly'' argument, is
Proposition~\ref{prop:lb}: \emph{for this concept class one always has
$m_{\mathrm{filt}}(D,\eps,\delta)\ge1$ whenever a target exists.} This is what
makes a clean multiplicative bound possible without adaptively shrinking the
labeled budget. The structure is: Section~\ref{sec:alg} gives the algorithm and
its labeled-sample upper bound (Lemma~\ref{lem:upper}); Section~\ref{sec:lower}
proves the structural lower bound $m_{\mathrm{filt}}\ge1$
(Proposition~\ref{prop:lb}) and assembles the multiplicative comparison
(Lemma~\ref{lem:compare}); Section~\ref{sec:thm} assembles the theorem.

\section{The algorithm and its labeled-sample upper bound}\label{sec:alg}

The classifier produced by a two-sided disjunction depends only on $h^+$, which
is a monotone disjunction. We therefore learn $h^{\ast+}$ directly with the
classical elimination (list-and-cross-off) algorithm. Correctness needs only
$h^\ast\in C$; the quantitative \emph{matching} of $m_{\mathrm{filt}}$ is
provided by Section~\ref{sec:lower}.

\begin{algorithm}[h]
\caption{Elimination learner for the classifier of $C$}\label{alg:elim}
\begin{algorithmic}[1]
  \Require parameters $\eps,\delta\in(0,1)$; oracle access to labeled examples
           $(x,y)$, $x\sim D$, $y=h^{\ast+}(x)$.
  \Ensure a classifier $g$ which is a monotone disjunction over $\{1,\dots,n\}$.
  \State $m \gets \bigl\lceil \tfrac1\eps\bigl(n\ln 2+\ln(1/\delta)\bigr)\bigr\rceil$
  \State $T \gets \{1,2,\dots,n\}$
         \Comment{$T$ = set of variables still allowed in $g$}
  \For{$t=1$ \textbf{to} $m$}
     \State draw a labeled example $(x,y)$
     \If{$y=-1$} \Comment{$h^{\ast+}(x)=0$}
        \State $T \gets T \setminus P(x)$
            \Comment{remove every variable set to $1$ in $x$}
     \EndIf
  \EndFor
  \State \Return the classifier $g(x):=\OR_{i\in T} x_i$
         \big(equivalently $g(x)=\mathbf 1[\,T\cap P(x)\neq\emptyset\,]$\big)
\end{algorithmic}
\end{algorithm}

\begin{lemma}\label{lem:upper}
Fix any $D$ on $X$ and any target $h^\ast\in C$. On a labeled sample of size
$m=\lceil\frac1\eps(n\ln2+\ln(1/\delta))\rceil$, Algorithm~\ref{alg:elim}:
\begin{enumerate}
\item runs in time $O(mn)=\mathrm{poly}(n,1/\eps,\log(1/\delta))$ and uses $m$
labeled examples and no unlabeled examples;
\item outputs a classifier $g$ with
$\Pr[\err_{D,h^\ast}(g)\le\eps]\ge 1-\delta$,
the probability over the labeled sample.
\end{enumerate}
\end{lemma}

\begin{proof}
\emph{Running time and sample size.} The loop runs $m$ times; each iteration
draws one labeled example, inspects the at most $n$ coordinates of $x$ and
deletes the corresponding elements from $T$, costing $O(n)$. Output is $O(n)$.
Total time $O(mn)$, polynomial in $n,1/\eps,\log(1/\delta)$, hence in
$n,1/\eps,1/\delta$. No unlabeled example is requested.

\emph{Structure of the output.} Let $S^\ast:=S^{\ast+}$ be the index set of the
target classifier $h^{\ast+}=\OR_{i\in S^\ast}x_i$, and let $T_{\mathrm{out}}$ be
the final value of $T$.

\textbf{(I1) $S^\ast\subseteq T_{\mathrm{out}}$.}
A variable $i$ is deleted only when, on some example $(x,-1)$, $i\in P(x)$. If
$i\in S^\ast$ then $x_i=1$ forces $h^{\ast+}(x)=1$, so the label would be $+1$,
not $-1$. Hence no $i\in S^\ast$ is ever deleted.

\textbf{(I2) $g$ is consistent with the whole sample.} For a sampled $(x,y)$:
if $y=+1$ then $S^\ast\cap P(x)\neq\emptyset$, and (I1) gives
$T_{\mathrm{out}}\cap P(x)\neq\emptyset$, so $g(x)=1$; if $y=-1$, the iteration
processing $(x,-1)$ deleted all of $P(x)$ and elements are never re-added, so
$T_{\mathrm{out}}\cap P(x)=\emptyset$ and $g(x)=0$. Thus $g$ agrees with the
target on every sampled example.

\emph{Generalization (finite-class realizable PAC bound).}
Let $\mathcal H=\{g_T:T\subseteq\{1,\dots,n\}\}$, $|\mathcal H|=2^n$; the output
$g\in\mathcal H$ and the target classifier $g_{S^\ast}\in\mathcal H$, so the
problem is realizable. Call $g_T$ \emph{bad} if $\err_{D,h^\ast}(g_T)>\eps$. For
a fixed bad $g_T$, $\Pr[g_T\text{ consistent with the sample}]
=(1-\err_{D,h^\ast}(g_T))^m<(1-\eps)^m\le e^{-\eps m}$. By a union bound over the
$\le2^n$ bad hypotheses,
$\Pr[\exists\text{ bad consistent }g_T]\le 2^n e^{-\eps m}$. With
$m\ge\frac1\eps(n\ln2+\ln(1/\delta))$ we get $\eps m\ge n\ln2+\ln(1/\delta)$,
hence $2^ne^{-\eps m}\le\delta$. By (I2) the output $g$ is consistent with the
sample, so with probability $\ge1-\delta$ it is not bad, i.e.\
$\err_{D,h^\ast}(g)\le\eps$.
\end{proof}

\section{The filtered complexity is at least one, and the multiplicative
comparison}\label{sec:lower}

The previous version of this argument compared $m$ to $m_{\mathrm{filt}}+1$,
which is vacuous: $m$ is already bounded by a polynomial $p$, so $m\le
p\le p\cdot(m_{\mathrm{filt}}+1)$ holds for \emph{any} learner with labeled usage
$\le p$, regardless of $m_{\mathrm{filt}}$. The additive ``$+1$'' was inserted to
cover a putative degenerate regime $m_{\mathrm{filt}}=0$. We now prove that this
regime \emph{cannot occur}, which removes the slack and turns the comparison into
a genuine multiplicative bound.

We begin with two structural facts about compatibility for this class.

\begin{lemma}[$0^n$ is incompatible]\label{lem:zero}
If $C_D\neq\emptyset$ then $0^n\notin\supp(D)$.
\end{lemma}

\begin{proof}
At $x=0^n$ we have $P(x)=\emptyset$, so $S^+\cap P(x)=S^-\cap P(x)=\emptyset$ for
\emph{every} pair $(S^+,S^-)$; thus $h^+(0^n)=h^-(0^n)=0$, violating the
compatibility requirement that exactly one of $h^+(x),h^-(x)$ equal $1$ at every
support point. Hence no hypothesis is compatible with a distribution whose
support contains $0^n$; if $C_D\neq\emptyset$ then $0^n\notin\supp(D)$.
\end{proof}

\begin{lemma}[Two extreme compatible targets]\label{lem:extreme}
Suppose $C_D\neq\emptyset$. Then both of the following hypotheses lie in $C_D$:
\[
  t_0:=(S^+=\emptyset,\;S^-=\{1,\dots,n\}),\qquad
  t_1:=(S^+=\{1,\dots,n\},\;S^-=\emptyset).
\]
Their classifiers $c_0:=t_0^{+}$ and $c_1:=t_1^{+}$ satisfy $c_0(x)=0$ and
$c_1(x)=1$ for every $x\in\supp(D)$; in particular $\err_{D,c_1}(c_0)=1$.
\end{lemma}

\begin{proof}
By Lemma~\ref{lem:zero}, $0^n\notin\supp(D)$, so every $x\in\supp(D)$ has
$P(x)\neq\emptyset$.

\emph{$t_0\in C_D$.} Here $t_0^+=\OR_{\emptyset}\equiv0$ and
$t_0^-=\OR_{\{1,\dots,n\}}$. For $x\in\supp(D)$: $t_0^+(x)=0$, and
$\{1,\dots,n\}\cap P(x)=P(x)\neq\emptyset$ so $t_0^-(x)=1$. Thus exactly one of
$t_0^+(x),t_0^-(x)$ is $1$, and $t_0$ is compatible. Its classifier is $c_0\equiv0$
on $\supp(D)$.

\emph{$t_1\in C_D$.} Here $t_1^+=\OR_{\{1,\dots,n\}}$ and $t_1^-\equiv0$. For
$x\in\supp(D)$: $t_1^+(x)=1$ (since $P(x)\neq\emptyset$) and $t_1^-(x)=0$. Again
exactly one is $1$, so $t_1$ is compatible, with classifier $c_1\equiv1$ on
$\supp(D)$.

Finally $\err_{D,c_1}(c_0)=\Pr_{x\sim D}[c_0(x)\neq c_1(x)]=\Pr_{x\sim
D}[0\neq1]=1$, since $c_0,c_1$ disagree at every support point.
\end{proof}

\begin{proposition}[The filtered complexity is at least one]\label{prop:lb}
For every $D$ with $C_D\neq\emptyset$ and all $\eps,\delta\in(0,\tfrac12)$,
\[
  m_{\mathrm{filt}}(D,\eps,\delta)\ge 1 .
\]
\end{proposition}

\begin{proof}
Suppose, for contradiction, that $m_{\mathrm{filt}}(D,\eps,\delta)=0$. By
Definition~\ref{def:filt} there is a randomized learning rule that uses $0$
labeled examples and outputs a classifier $g$ (a random function, whose
distribution does not depend on the unknown target, since it sees no
target-dependent input) such that for \emph{every} target $h^\ast\in C_D$,
$\Pr[\err_{D,h^\ast}(g)\le\eps]\ge1-\delta>0$; in particular for each target
there is at least one realization $g$ of the output with
$\err_{D,h^\ast}(g)\le\eps$.

Apply this to the two targets $t_0,t_1\in C_D$ of Lemma~\ref{lem:extreme}, with
classifiers $c_0\equiv0$, $c_1\equiv1$ on $\supp(D)$. For \emph{any} fixed
classifier $g$,
\[
  \err_{D,t_0}(g)+\err_{D,t_1}(g)
  =\Pr_{x\sim D}[g(x)\neq0]+\Pr_{x\sim D}[g(x)\neq1]
  =\Pr_{x\sim D}[g(x)=1]+\Pr_{x\sim D}[g(x)=0]
  =1,
\]
where the middle equality uses $c_0(x)=0,c_1(x)=1$ on $\supp(D)$ and the last
that the two events partition $\supp(D)$. Hence, for every realized $g$,
$\max\{\err_{D,t_0}(g),\err_{D,t_1}(g)\}\ge\tfrac12>\eps$ (as $\eps<\tfrac12$).
Therefore \emph{no} realization of the output can satisfy both
$\err_{D,t_0}(g)\le\eps$ and $\err_{D,t_1}(g)\le\eps$. Writing $A_j:=\{g:\err_{D,t_j}(g)\le\eps\}$, the displayed identity
shows that the sets $A_0,A_1$ of \emph{successful} outputs are disjoint:
$A_0\cap A_1=\emptyset$ (no $g$ lies in both, since that would give
$\err_{D,t_0}(g)+\err_{D,t_1}(g)\le2\eps<1$, contradicting the identity
$=1$). Since the learner's output distribution is identical against $t_0$
and against $t_1$ (it receives no labels), $\Pr[A_0]$ and $\Pr[A_1]$ are
well defined under this common distribution, and
$\Pr[A_0]+\Pr[A_1]=\Pr[A_0\cup A_1]\le1$.

Since $\delta<\tfrac12$ we have $1-\delta>\tfrac12$, so the success
guarantee demands $\Pr[A_0]\ge1-\delta>\tfrac12$ and
$\Pr[A_1]\ge1-\delta>\tfrac12$ simultaneously, whence
$\Pr[A_0]+\Pr[A_1]>1$. This contradicts $\Pr[A_0]+\Pr[A_1]\le1$
(from $A_0\cap A_1=\emptyset$). Hence the assumption
$m_{\mathrm{filt}}(D,\eps,\delta)=0$ is impossible, and
$m_{\mathrm{filt}}(D,\eps,\delta)\ge1$.
\end{proof}

\begin{remark}
The clean takeaway of Proposition~\ref{prop:lb}, used below, is:
\emph{for every $D$ with $C_D\neq\emptyset$ and every $\eps,\delta\in(0,\tfrac12)$,
one has $m_{\mathrm{filt}}(D,\eps,\delta)\ge1$.} This is the only consequence of
the proposition that the comparison invokes, and the proof of this consequence is
exactly the contradiction with $A_0\cap A_1=\emptyset$ above for $\delta<\tfrac12$.
(In fact a standard two-point/EHKV refinement using the everywhere-disagreeing
pair $t_0,t_1$ gives the stronger
$m_{\mathrm{filt}}(D,\eps,\delta)=\Omega\bigl(\tfrac1\eps\log\tfrac1\delta\bigr)$,
but $m_{\mathrm{filt}}\ge1$ already suffices for the multiplicative bound.)
\end{remark}

\begin{lemma}[Multiplicative comparison]\label{lem:compare}
Let $p(n,1/\eps,\log(1/\delta)):=\frac1\eps\bigl(n+\log_2(1/\delta)+1\bigr)$.
For every $D$ with $C_D\neq\emptyset$ and all $\eps,\delta\in(0,\tfrac12)$,
\[
  m=\Bigl\lceil \tfrac1\eps\bigl(n\ln 2+\ln(1/\delta)\bigr)\Bigr\rceil
  \;\le\; p\bigl(n,1/\eps,\log(1/\delta)\bigr)\cdot
          m_{\mathrm{filt}}(D,\eps,\delta).
\]
\end{lemma}

\begin{proof}
Using $\lceil a\rceil\le a+1$, $\ln2<1$, $\ln(1/\delta)<\log_2(1/\delta)$, and
$1\le\frac1\eps$ (as $\eps<1$),
\[
  m \le \tfrac1\eps\bigl(n\ln2+\ln(1/\delta)\bigr)+1
     \le \tfrac1\eps\bigl(n+\log_2(1/\delta)\bigr)+\tfrac1\eps
     = \tfrac1\eps\bigl(n+\log_2(1/\delta)+1\bigr)
     = p\bigl(n,1/\eps,\log(1/\delta)\bigr).
\]
By Proposition~\ref{prop:lb}, $m_{\mathrm{filt}}(D,\eps,\delta)\ge1$. Therefore
\[
  m \le p\bigl(n,1/\eps,\log(1/\delta)\bigr)
     = p\bigl(n,1/\eps,\log(1/\delta)\bigr)\cdot 1
     \le p\bigl(n,1/\eps,\log(1/\delta)\bigr)\cdot
         m_{\mathrm{filt}}(D,\eps,\delta).\qedhere
\]
\end{proof}

The point is that $m_{\mathrm{filt}}\ge1$ is exactly what licenses replacing
$(m_{\mathrm{filt}}+1)$ by $m_{\mathrm{filt}}$: the comparison is now genuinely
multiplicative, with no additive slack, and $m_{\mathrm{filt}}$ appears
non-trivially as a multiplicative lower-anchor on the right-hand side.

\section{Proof of Theorem~\ref{thm:main}}\label{sec:thm}

Take $\mathcal A$ to be Algorithm~\ref{alg:elim} with parameters
$\eps,\delta\in(0,\tfrac12)$. Fix any $D$ with $C_D\neq\emptyset$ and any target
$h^\ast\in C_D\subseteq C$.

\emph{Clause (a).} By Lemma~\ref{lem:upper}(1) the running time is $O(mn)$ with
$m=\lceil\frac1\eps(n\ln2+\ln(1/\delta))\rceil$, polynomial in $n,1/\eps$ and
(since $\log(1/\delta)\le1/\delta$) in $1/\delta$.

\emph{Correctness and clause (b).} By Lemma~\ref{lem:upper}(2), $\mathcal A$
outputs a classifier $g$ with $\Pr[\err_{D,h^\ast}(g)\le\eps]\ge1-\delta$, using
exactly $m$ labeled examples and no unlabeled examples (Lemma~\ref{lem:upper}
needs only $h^\ast\in C$, which holds since $C_D\subseteq C$). By
Lemma~\ref{lem:compare},
\[
  m \le p\bigl(n,1/\eps,\log(1/\delta)\bigr)\cdot
        m_{\mathrm{filt}}(D,\eps,\delta),
  \qquad p\bigl(n,1/\eps,\log(1/\delta)\bigr)=\tfrac1\eps\bigl(n+\log_2(1/\delta)+1\bigr),
\]
a fixed polynomial in $n,1/\eps,\log(1/\delta)$. Hence the labeled-sample size
matches the compatibility-filtered sample complexity up to a polynomial factor in
$n,1/\eps,\log(1/\delta)$, in the genuine multiplicative sense.

\emph{Unlabeled sample.} $\mathcal A$ uses $0$ unlabeled examples, trivially a
polynomial number.

If $C_D=\emptyset$ there is no compatible target and the learning instance is
empty, so there is nothing to prove. In all instances with a target, all three
requirements hold, completing the proof of Theorem~\ref{thm:main}. \qed

\begin{remark}[Why no adaptivity or unlabeled data is needed --- and why this is
not vacuous]
A priori, matching $m_{\mathrm{filt}}$ \emph{multiplicatively} could require an
adaptive learner that detects when compatibility collapses the live class
$C_D$ to a single classifier (the regime $m_{\mathrm{filt}}=0$) and then spends
no labeled examples; otherwise a learner that always draws $\Theta(n/\eps)$
labeled examples would violate $m_{\mathrm A}\le p\cdot m_{\mathrm{filt}}$ at
$m_{\mathrm{filt}}=0$. Proposition~\ref{prop:lb} shows that for two-sided
disjunctions this collapse \emph{never happens}: compatibility forces the two
everywhere-disagreeing targets $t_0,t_1$ into $C_D$ (Lemma~\ref{lem:extreme}),
so $C_D$ always contains a maximally separated pair and $m_{\mathrm{filt}}\ge1$.
Consequently the non-adaptive, unlabeled-data-free elimination learner already
attains a true multiplicative match. The semantic content is thus that, for this
particular class, the compatibility structure does not in fact shrink the
classifier complexity below the universal floor $m_{\mathrm{filt}}\ge1$, so
unlabeled data buys nothing here --- in contrast to general Balcan--Blum settings
where it can. This is a statement about the class, established by
Proposition~\ref{prop:lb}, not an artifact of an additive ``$+1$''.
\end{remark}

\end{document}
